\ifdefined\gkver\else\def\gkver{student}\fi
\documentclass[\gkver]{gaokaozhenti}
\begin{document}

\chapter{2009年北京理综}
%% paperType: 理综物理部分

\begin{enumerate}
\item
%% number: 13
%% paperName: 2009年北京理综第 13 题 4 分
%% typeId: 1
%% type: 单选题
%% chapter: 气体性质
%% point: 分子动理论
%% method: 无
%% score: 4
%% degree: 650
%% duplicateId: 0
%% body:
做布朗运动实验，得到某个观察记录如图，图中记录的是\xzanswer{D}

\onepicture[w=4.6cm]{figs/13.png}

\fourchoices[answer=D]
{分子无规则运动的情况}
{某个微粒做布朗运动的轨迹}
{某个微粒做布朗运动的速度-时间图线}
{按等时间间隔依次记录的某个运动微粒位置的连线}

%% answer: D
%% memo:
\memoanswer{
布朗运动是悬浮在液体中的固体小颗粒的无规则运动，而非分子的运动，故 A 项错误；既然运动是无规则的，所以微粒没有固定的运动轨迹，故 B 项错误；对于某个微粒而言在不同时刻的速度大小和方向均是不确定的，所以无法确定其在某一个时刻的速度，故也就无法描绘其速度—时间图线，故 C 项错误；只有 D 项正确。
}

\item
%% number: 14
%% paperName: 2009年北京理综第 14 题 4 分
%% typeId: 1
%% type: 单选题
%% chapter: 原子和原子核
%% point: 天然放射性
%% method: 无
%% score: 4
%% degree: 650
%% duplicateId: 0
%% body:
下列现象中，与原子核内部变化有关的是\xzanswer{B}

\fourchoices[answer=B]
{$\alpha$ 粒子散射现象}
{天然放射现象}
{光电效应现象}
{原子发光现象}

%% answer: B
%% memo:
\memoanswer{
$\alpha$ 粒子散射实验表明了原子内部有一个很小的核，并没有涉及核内部的变化，故 A 项错误；天然放射现象是原子核内部发生变化自发地放射出 $\alpha$ 粒子或电子，从而发生 $\alpha$ 衰变或 $\beta$ 衰变，故 B 项正确；光电效应是原子核外层电子获得光子能量脱离原子核的束缚而逸出，没有涉及原子核的变化，故 C 项错误；原子发光是原子跃迁形成的，也没有涉及原子核的变化，故 D 项错误。
}

\item
%% number: 15
%% paperName: 2009年北京理综第 15 题 4 分
%% typeId: 1
%% type: 单选题
%% chapter: 光学
%% point: 光的电磁说
%% method: 类比
%% score: 4
%% degree: 650
%% duplicateId: 0
%% body:
类比是一种有效的学习方法，通过归纳和比较，有助于掌握新知识，提高学习效率。在类比过程中，既要找出共同之处，又要抓住不同之处。某同学对机械波和电磁波进行类比，总结出下列内容，其中不正确的是\xzanswer{D}

\fourchoices[answer=D]
{机械波的频率、波长和波速三者满足的关系，对电磁波也适用}
{机械波和电磁波都能产生干涉和衍射现象}
{机械波的传播依赖于介质，而电磁波可以在真空中传播}
{机械波既有横波又有纵波，而电磁波只有纵波}

%% answer: D
%% memo:
\memoanswer{
波长、波速、频率的关系对任何波都是成立的，对电磁波当然成立，故 A 项正确；干涉和衍射是波的特性，机械波、电磁波都是波，这些特性都具有，故 B 项正确；机械波是机械振动在介质中传播形成的，所以机械波的传播需要介质，而电磁波是交替变化的电场和磁场由近及远的传播形成的，所以电磁波传播不需要介质，故 C 项正确；机械波既有横波又有纵波，但是电磁波只能是横波，其证据就是电磁波能够发生偏振现象，而偏振现象是横波才有的，D 项错误。故正确答案应为 D。
}

\item
%% number: 16
%% paperName: 2009年北京理综第 16 题 4 分
%% typeId: 1
%% type: 单选题
%% chapter: 电场
%% point: 电场线
%% method: 无
%% score: 4
%% degree: 650
%% duplicateId: 0
%% body:
某静电场的电场线分布如图所示，图中 P、Q 两点的电场强度分别为 $E_{P}$ 和 $E_{Q}$，电势分别为 $U_{P}$ 和 $U_{Q}$。则\xzanswer{A}

\onepicture[w=4.8cm]{figs/16.png}

\fourchoices[answer=A]
{$E_{P}>E_{Q}$，$U_{P}>U_{Q}$}
{$E_{P}>E_{Q}$，$U_{P}<U_{Q}$}
{$E_{P}<E_{Q}$，$U_{P}>U_{Q}$}
{$E_{P}<E_{Q}$，$U_{P}<U_{Q}$}

%% answer: A
%% memo:
\memoanswer{
从图可以看出 P 点的电场线的密集程度大于 Q 点的密集程度，故 P 点的场强大于 Q 点的场强，因电场线的方向由 P 指向 Q，而沿电场线的方向电势逐渐降低，P 点的电势高于 Q 点的电势，故 A 项正确。
}

\item
%% number: 17
%% paperName: 2009年北京理综第 17 题 4 分
%% typeId: 1
%% type: 单选题
%% chapter: 机械振动
%% point: 振动图象
%% method: 无
%% score: 4
%% degree: 650
%% duplicateId: 0
%% body:
一简谐机械波沿 $x$ 轴正方向传播，周期为 $T$，波长为 $\lambda$，若在 $x=0$ 处质点的振动图像如右图所示，则该波在 $t=T/2$ 时的波形曲线为\xzanswer{A}

\onepicture[w=5.2cm]{figs/17a.png}

\fourchoices[answer=A, ispicture=true, h=2.4cm]
{figs/17b.png}
{figs/17c.png}
{figs/17d.png}
{figs/17e.png}

%% answer: A
%% memo:
\memoanswer{
从振动图上可以看出 $x=0$ 处的质点在 $t=T/2$ 时刻处于平衡位置，且正在向下振动，四个选项中只有 A 图符合要求，故 A 项正确。
}

\item
%% number: 18
%% paperName: 2009年北京理综第 18 题 4 分
%% typeId: 1
%% type: 单选题
%% chapter: 运动和力
%% point: 牛顿第二定律
%% method: 无
%% score: 4
%% degree: 650
%% duplicateId: 0
%% body:
如图所示，将质量为 $m$ 的滑块放在倾角为 $\theta$ 的固定斜面上，滑块与斜面之间的动摩擦因数为 $\mu$，若滑块与斜面之间的最大静摩擦力与滑动摩擦力大小相等，重力加速度为 $g$，则\xzanswer{C}

\onepicture[w=4.6cm]{figs/18.png}

\fourchoices[answer=C]
{将滑块由静止释放，如果 $\mu>\tan\theta$，滑块将下滑}
{给滑块沿斜面向下的初速度，如果 $\mu<\tan\theta$，滑块将减速下滑}
{用平行于斜面向上的力拉滑块向上匀速运动，如果 $\mu=\tan\theta$，拉力大小应是 $2mg\sin\theta$}
{用平行于斜面向下的力拉滑块向下匀速运动，如果 $\mu=\tan\theta$，拉力大小应是 $mg\sin\theta$}

%% answer: C
%% memo:
\memoanswer{
对处于斜面上的物块受力分析，要使物块沿斜面下滑则 $mg\sin\theta>\mu mg\cos\theta$，故 $\mu<\tan\theta$，故 A、B 错误；若要使物块在平行于斜面向上的拉力 $F$ 的作用下沿斜面匀速上滑，由平衡条件有 $F-mg\sin\theta-\mu mg\cos\theta=0$，故 $F=mg\sin\theta+\mu mg\cos\theta$，若 $\mu=\tan\theta$，则 $mg\sin\theta=\mu mg\cos\theta$，即 $F=2mg\sin\theta$，故 C 项正确；若要使物块在平行于斜面向下的拉力 $F$ 作用下沿斜面向下匀速滑动，由平衡条件有 $F+mg\sin\theta-\mu mg\cos\theta=0$，则 $F=\mu mg\cos\theta-mg\sin\theta$，若 $\mu=\tan\theta$，则 $mg\sin\theta=\mu mg\cos\theta$，即 $F=0$，故 D 项错误。
}

\item
%% number: 19
%% paperName: 2009年北京理综第 19 题 4 分
%% typeId: 1
%% type: 单选题
%% chapter: 磁场
%% point: 带电粒子在复合场中的运动
%% method: 无
%% score: 4
%% degree: 650
%% duplicateId: 0
%% body:
如图所示的虚线区域内，充满垂直于纸面向里的匀强磁场和竖直向下的匀强电场，一带电粒子 a（不计重力）以一定的初速度由左边界的 O 点射入磁场、电场区域，恰好沿直线由区域右边界的 $O'$ 点（图中未标出）穿出。若撤去该区域内的磁场，而保留电场不变，另一个同样的粒子 b（不计重力）仍以相同初速度由 O 点射入，从区域右边界穿出，则粒子 b\xzanswer{C}

\onepicture[w=4.4cm]{figs/19.png}

\fourchoices[answer=C]
{穿出位置一定在 $O'$ 下方}
{穿出位置一定在 $O'$ 上方}
{运动时，在电场中的电势能一定减小}
{在电场中运动时，动能一定减小}

%% answer: C
%% memo:
\memoanswer{
a 粒子要在电场、磁场的复合场区内做直线运动，则该粒子一定做匀速直线运动，故对粒子 a 有 $Bqv=Eq$，即只要满足 $E=Bv$，无论粒子带正电还是负电，粒子都可以沿直线穿出复合场区；当撤去磁场只保留电场时，粒子 b 由于电性不确定，故无法判断从 $O'$ 点的上方或下方穿出，故 A、B 错误；粒子 b 在穿过电场区的过程中必然受到电场力的作用而做类平抛运动，电场力做正功，其电势能减小，动能增大，故 C 项正确、D 项错误。
}

\item
%% number: 20
%% paperName: 2009年北京理综第 20 题 4 分
%% typeId: 1
%% type: 单选题
%% chapter: 电场
%% point: 电场叠加
%% method: 无
%% score: 4
%% degree: 350
%% duplicateId: 0
%% body:
如图所示为一个内、外半径分别为 $R_{1}$ 和 $R_{2}$ 的圆环状均匀带电平面，其单位面积带电量为 $\sigma$。取环面中心 O 为坐标原点，以垂直于环面的轴线为 $x$ 轴。设轴上任意点 P 到 O 的距离为 $x$，P 点电场强度的大小为 $E$。下面给出 $E$ 的四个表达式（式中 $k$ 为静电力常量），其中只有一个是合理的。你可能不会求解此处的场强 $E$，但是你可以通过一定的物理分析，对下列表达式的合理性做出判断。根据你的判断，$E$ 的合理表达式应为\xzanswer{B}

\onepicture[w=5.0cm]{figs/20.png}

\fourchoices[answer=B]
{$E=2\pi\sigma k\left(\frac{R_{1}}{\sqrt{x^{2}+R_{1}^{2}}}-\frac{R_{2}}{\sqrt{x^{2}+R_{2}^{2}}}\right)x$}
{$E=2\pi\sigma k\left(\frac{1}{\sqrt{x^{2}+R_{1}^{2}}}-\frac{1}{\sqrt{x^{2}+R_{2}^{2}}}\right)x$}
{$E=2\pi\sigma k\left(\frac{R_{1}}{\sqrt{x^{2}+R_{1}^{2}}}+\frac{R_{2}}{\sqrt{x^{2}+R_{2}^{2}}}\right)x$}
{$E=2\pi\sigma k\left(\frac{1}{\sqrt{x^{2}+R_{1}^{2}}}+\frac{1}{\sqrt{x^{2}+R_{2}^{2}}}\right)x$}

%% answer: B
%% memo:
\memoanswer{
当 $R_{1}=0$ 时，对于 A 项而言 $E=0$，此时带电圆环演变为带电圆面，中心轴线上一点的电场强度 $E>0$，故 A 项错误；当 $x=0$ 时，此时要求的场强为 O 点的场强，由对称性可知 $E_{0}=0$，对于 C 项而言，$x=0$ 时 $E$ 为一定值，故 C 项错误；当 $x\to\infty$ 时 $E\to0$，而 D 项中 $E\to4\pi k\sigma$，故 D 项错误；所以正确选项只能为 B。
}

\item
%% number: 21
%% paperName: 2009年北京理综第 21 题 12 分
%% typeId: 4
%% type: 实验题
%% chapter: 电路
%% point: 测电动势与内阻
%% method: 无
%% score: 12
%% degree: 450
%% duplicateId: 0
%% body:
\begin{enumerate}
	\item
	Ⅰ．（4 分）在《用双缝干涉测量光的波长》实验中，将双缝干涉仪按要求安装在光具座上（如图 1），并选用缝间距 $d=0.2\Umm$ 的双缝屏。从仪器注明的规格可知，像屏与双缝屏间的距离 $L=700\Umm$。然后接通电源使光源正常工作。

	\onepicture[w=11.0cm]{figs/21a.png}

	\begin{steps}
		\item
		已知测量头主尺的最小刻度是毫米，副尺上有 50 个分度，某同学调整手轮后，从测量头的目镜看去，第一次映入眼帘的干涉条纹如图 2（a）所示，图 2（a）中的数字是该同学给各暗纹的编号，此时图 2（b）中的游标尺上的读数 $x_{1}=1.16\Umm$，接着再转动手轮，映入眼帘的干涉条纹如图 3（a）所示，此时图 3（b）中的游标尺上的读数 $x_{2}=$\tkanswer[1.6cm]{} $\Umm$。
		\item
		利用上述测量结果，经计算可得两个相邻明纹（或暗纹）间的距离 $\Delta x=$\tkanswer[1.6cm]{} $\Umm$，这种色光的波长 $\lambda=$\tkanswer[1.6cm]{} $\Unm$。
	\end{steps}

	\twopicture[showlabel=false, w=7.4cm]
	{figs/21b.png}
	{figs/21c.png}

	\item
	Ⅱ．（12 分）某同学通过查找资料自己动手制作了一个电池，该同学想测量一下这个电池的电动势 $E$ 和内电阻 $r$。但是从实验室只借到一个开关、一个电阻箱（最大电阻为 $999.9\UO$，可当标准电阻用）、一只电流表（量程 $I_{A}=0.6\UA$，内阻 $r_{A}=0.1\UO$）和若干导线。

	\begin{steps}
		\item
		请根据测量电动势 $E$ 和内电阻 $r$ 的要求，设计图 1 中器件的连接方式。画线把它们接起来。
		\item
		接通开关，逐次改变电阻箱的阻值 $R$，读出与 $R$ 对应的电流表的示数 $I$，并作记录。当电阻箱的阻值 $R=2.6\UO$ 时，其对应的电流表的示数如图 2 所示。处理实验数据时，首先计算出每个电流值 $I$ 的倒数 $\frac{1}{I}$，再制作 $R-\frac{1}{I}$ 坐标图，如图 3 所示，图中已标出了（$R$，$\frac{1}{I}$）的几个与测量对应的坐标点，请你将与图 2 实验数据对应的坐标点也标在图 3 上。
		\item
		在图 3 上把描绘出的坐标点连成图线。
		\item
		根据图 3 描绘出的图线，可以得出这个电池的电动势 $E=$\tkanswer[1.6cm]{} $\UV$；内电阻 $r=$\tkanswer[1.6cm]{} $\UO$。
	\end{steps}

	\threepicture[showlabel=false, w=6.4cm]
	{figs/21d.png}
	{figs/21e.png}
	{figs/21f.png}
\end{enumerate}

%% answer:
\jdanswer{
\begin{enumerate}
	\item
	①15.02；②2.31，$6.6\times10^{2}$
	\item
	①图略；②图略；③图略；④$E=1.5\UV$，$r=0.3\UO$。
\end{enumerate}
}

%% memo:
\memoanswer{
（1）由游标卡尺的读数规则可知 $x_{2}=15.0\Umm+1\times0.02\Umm=15.02\Umm$；图 2（a）中暗纹与图 3（a）中暗纹间的间隔为 6 个，故 $\Delta x=(x_{2}-x_{1})/6=(15.02-1.16)/6=2.31\Umm$；由 $\Delta x=L\lambda/d$ 可知 $\lambda=d\Delta x/L=0.20\Umm\times2.31\Umm/700\Umm=6.6\times10^{2}\Unm$。

（2）根据闭合电路欧姆定律，测量电源的电动势和内电阻，需要得到电源的路端电压和通过电源的电流，在本实验中没有电压表，但是可以用电阻箱和电流表串联充当电压表，测量电源的路端电压，通过电流表的电流也是通过电源的电流，所以只需要将电流表和电阻箱串联接在电源两端即可。实物图的连接如答图所示。

由闭合电路欧姆定律有 $E=I(R+r+r_{g})$，解得 $R=E\cdot\frac{1}{I}-(r+r_{g})$，根据 $R-\frac{1}{I}$ 图线可知：电源的电动势等于图线的斜率，内阻为纵轴负方向的截距减去电流表的内阻。

\onepicture[w=5.6cm]{figs/21_an1.jpg}

\onepicture[w=5.6cm]{figs/21_an2.jpg}
}

\item
%% number: 22
%% paperName: 2009年北京理综第 22 题 16 分
%% typeId: 6
%% type: 计算题
%% chapter: 曲线运动
%% point: 人造地球卫星
%% method: 无
%% score: 16
%% degree: 450
%% duplicateId: 0
%% body:
已知地球半径为 $R$，地球表面重力加速度为 $g$，不考虑地球自转的影响。

\begin{enumerate}
	\item
	推导第一宇宙速度 $v_{1}$ 的表达式；
	\item
	若卫星绕地球做匀速圆周运动，运行轨道距地面高度为 $h$，求卫星的运行周期 $T$。
\end{enumerate}

%% answer:
\jdanswer{
\begin{enumerate}
	\item
	$v_{1}=\sqrt{gR}$
	\item
	$T=\frac{2\pi}{R}\sqrt{\frac{(R+h)^{2}}{g}}$
\end{enumerate}
}

%% memo:
\memoanswer{
（1）设卫星的质量为 $m$，地球的质量为 $M$，在地球表面附近满足 $G\frac{Mm}{R^{2}}=mg$，得 $GM=R^{2}g$ ①。

卫星做圆周运动的向心力等于它受到的万有引力 $m\frac{v_{1}^{2}}{R}=G\frac{Mm}{R^{2}}$ ②。

①式代入②式，得到 $v_{1}=\sqrt{Rg}$。

（2）考虑①式，卫星受到的万有引力为 $F=G\frac{Mm}{(R+h)^{2}}=\frac{mgR^{2}}{(R+h)^{2}}$ ③。

由牛顿第二定律 $F=m\frac{4\pi^{2}}{T^{2}}(R+h)$ ④。

③、④联立解得 $T=\frac{2\pi}{R}\sqrt{\frac{(R+h)^{2}}{g}}$。
}

\item
%% number: 23
%% paperName: 2009年北京理综第 23 题 18 分
%% typeId: 6
%% type: 计算题
%% chapter: 电磁感应
%% point: 切割磁感线电动势
%% method: 无
%% score: 18
%% degree: 450
%% duplicateId: 0
%% body:
单位时间内流过管道横截面的液体体积叫做液体的体积流量（以下简称流量）。有一种利用电磁原理测量非磁性导电液体（如自来水、啤酒等）流量的装置，称为电磁流量计。它主要由将流量转换为电压信号的传感器和显示仪表两部分组成。

传感器的结构如图所示，圆筒形测量管内壁绝缘，其上装有一对电极 a 和 c，a、c 间的距离等于测量管的直径 $D$，测量管的轴线与 a、c 的连线方向以及通电线圈产生的磁场方向三者相互垂直。当导电液体流过测量管时，在电极 a、c 间出现感应电动势 $E$，并通过与电极连接的仪表显示出液体的流量 $Q$。设磁场均匀恒定，磁感应强度为 $B$。

\begin{enumerate}
	\item
	已知 $D=0.4\Um$，$B=2.5\times10^{-3}\UT$，$Q=0.12\Umc/\Us$，设液体在测量管内各处流速相同，试求 $E$ 的大小（$\pi$ 取 3.0）；
	\item
	一新建供水站安装了电磁流量计，在向用户供水时流量本应显示为正值，但实际显示却为负值。经检查，是误将测量管接反了，即液体从测量管出水口流入，入水口流出。因水已加压充满管道，不便再将测量管拆下重装，请你提出使显示仪表的流量示数变为正值的简便方法；
	\item
	显示仪表相当于传感器负载电阻，其阻值记为 $R$。a、c 间导电液体的电阻 $r$ 随液体电阻率的变化而变化，从而会影响显示仪表的示数。试以 $E$、$R$、$r$ 为参量给出电极 a、c 间输出电压 $U$ 的表达式，并说明怎样可以降低液体电阻率变化对显示仪表示数的影响？
\end{enumerate}

\onepicture[w=5.2cm, align=right]{figs/23.png}

%% answer:
\jdanswer{
\begin{enumerate}
	\item
	$E=1.0\times10^{-3}\UV$
	\item
	改变通电线圈中电流的方向，使磁场 $B$ 反向，或将传感器输出端对调接入显示仪表。
	\item
	$U=\frac{E}{1+r/R}$，见解析。
\end{enumerate}
}

%% memo:
\memoanswer{
（1）导电液体通过测量管时，相当于导线做切割磁感线的运动，在电极 a、c 间切割磁感线的液柱长度为 $D$，设液体的流速为 $v$，则产生的感应电动势为 $E=BDv$ ①。

由流量的定义，有 $Q=Sv=\frac{\pi D^{2}}{4}v$ ②。

①②式联立解得 $E=BD\frac{4Q}{\pi D^{2}}=\frac{4BQ}{\pi D}$。

代入数据得 $E=\frac{4\times2.5\times10^{-3}\times0.12}{3\times0.4}\UV=1.0\times10^{-3}\UV$。

（2）能使仪表显示的流量变为正值的简便方法（合理即可），如：改变通电线圈中电流的方向，使磁场 $B$ 反向，或将传感器输出端对调接入显示仪表。

（3）传感器的显示仪表构成闭合电路，由闭合电路欧姆定律有 $I=\frac{E}{R+r}$，$U=IR=\frac{RE}{R+r}=\frac{E}{1+r/R}$ ③。

输入显示仪表的是 a、c 间的电压 $U$，流量示数和 $U$ 一一对应，$E$ 与液体电阻率无关，而 $r$ 随电阻率的变化而变化，由③式可看出，$r$ 变化相应的 $U$ 也随之变化。在实际流量不变的情况下，仪表显示的流量示数会随 a、c 间的电压 $U$ 的变化而变化，增大 $R$，使 $R\gg r$，则 $U\approx E$，这样就可以降低液体电阻率的变化对显示仪表流量示数的影响。
}

\item
%% number: 24
%% paperName: 2009年北京理综第 24 题 20 分
%% typeId: 6
%% type: 计算题
%% chapter: 运动和力
%% point: 动量守恒定律
%% method: 无
%% score: 20
%% degree: 250
%% duplicateId: 0
%% body:
\begin{enumerate}
	\item
	如图 1 所示，ABC 为一固定在竖直平面内的光滑轨道，BC 段水平，AB 段与 BC 段平滑连接。质量为 $m_{1}$ 的小球从高为 $h$ 处由静止开始沿轨道下滑，与静止在轨道 BC 段上质量为 $m_{2}$ 的小球发生碰撞，碰撞前后两球的运动方向处于同一水平线上，且在碰撞过程中无机械能损失，求碰撞后小球 $m_{2}$ 的速度大小 $v_{2}$。
	\item
	碰撞过程中的能量传递规律在物理学中有着广泛的应用。为了探究这一规律，我们采用多球依次碰撞、碰撞前后速度在同一直线上、且无机械能损失的简化模型。如图 2 所示，在固定光滑水平直轨道上，质量分别为 $m_{1}$、$m_{2}$、……$m_{n-1}$、$m_{n}$、……的若干个球沿直线相间排列，给第一个球初动能 $E_{k1}$，从而引起各球的依次碰撞。定义其中第 $n$ 个球经过一次碰撞后获得的动能 $E_{kn}$ 与 $E_{k1}$ 之比为第一个球对第 $n$ 个球的动能传递系数 $k_{1n}$。

	a．求 $k_{1n}$；

	b．若 $m_{1}=4m_{0}$，$m_{3}=3m_{0}$，$m_{0}$ 为确定的已知量，求 $m_{2}$ 为何值时，$k_{13}$ 最大。
\end{enumerate}
\twopicture[showlabel=false, align=right, w=6.2cm]
{figs/24a.png}
{figs/24b.png}

%% answer:
\jdanswer{
\begin{enumerate}
	\item
	$v_{2}=\frac{2m_{1}\sqrt{2gh}}{m_{1}+m_{2}}$
	\item
	a．$k_{1n}=\frac{4^{n-1}m_{1}m_{2}^{2}m_{3}^{2}\cdots m_{n-1}^{2}m_{n}}{(m_{1}+m_{2})^{2}(m_{2}+m_{3})^{2}\cdots(m_{n-1}+m_{n})^{2}}$；

	b．$m_{2}=2m_{0}$ 时，$k_{13}$ 最大。
\end{enumerate}
}

%% memo:
\memoanswer{
（1）设碰撞前的速度为 $v_{10}$，根据机械能守恒定律 $m_{1}gh=\frac{1}{2}m_{1}v_{10}^{2}$ ①。

设碰撞后 $m_{1}$ 与 $m_{2}$ 的速度分别为 $v_{1}$ 和 $v_{2}$，根据动量守恒定律 $m_{1}v_{10}=m_{1}v_{1}+m_{2}v_{2}$ ②。

由于碰撞过程中无机械能损失，$\frac{1}{2}m_{1}v_{10}^{2}=\frac{1}{2}m_{1}v_{1}^{2}+\frac{1}{2}m_{2}v_{2}^{2}$ ③。

②、③式联立解得 $v_{2}=\frac{2m_{1}v_{10}}{m_{1}+m_{2}}$ ④，将①代入④得 $v_{2}=\frac{2m_{1}\sqrt{2gh}}{m_{1}+m_{2}}$。

（2）a．由④式，考虑到 $E_{k1}=\frac{1}{2}m_{1}v_{10}^{2}$ 和 $E_{k2}=\frac{1}{2}m_{2}v_{2}^{2}$，根据动能传递系数的定义，对于 1、2 两球 $k_{12}=\frac{E_{k2}}{E_{k1}}=\frac{4m_{1}m_{2}}{(m_{1}+m_{2})^{2}}$ ⑤。

同理可得，球 $m_{2}$ 和球 $m_{3}$ 碰撞后，动能传递系数 $k_{13}$ 应为 $k_{13}=\frac{E_{k3}}{E_{k1}}=\frac{E_{k2}}{E_{k1}}\cdot\frac{E_{k3}}{E_{k2}}=\frac{4m_{1}m_{2}}{(m_{1}+m_{2})^{2}}\cdot\frac{4m_{2}m_{3}}{(m_{2}+m_{3})^{2}}$ ⑥。

依次类推，动能传递系数 $k_{1n}$ 应为 $k_{1n}=\frac{E_{kn}}{E_{k1}}=\frac{E_{k2}}{E_{k1}}\cdot\frac{E_{k3}}{E_{k2}}\cdots\frac{E_{kn}}{E_{k(n-1)}}=\frac{4m_{1}m_{2}}{(m_{1}+m_{2})^{2}}\cdot\frac{4m_{2}m_{3}}{(m_{2}+m_{3})^{2}}\cdots\frac{4m_{n-1}m_{n}}{(m_{n-1}+m_{n})^{2}}$。

解得 $k_{1n}=\frac{4^{n-1}m_{1}m_{2}^{2}m_{3}^{2}\cdots m_{n-1}^{2}m_{n}}{(m_{1}+m_{2})^{2}(m_{2}+m_{3})^{2}\cdots(m_{n-1}+m_{n})^{2}}$。

b．将 $m_{1}=4m_{0}$、$m_{3}=m_{0}$ 代入⑥式可得 $k_{13}=64m_{0}^{2}\left[\frac{m_{2}}{(4m_{0}+m_{2})(m_{2}+m_{0})}\right]^{2}$。

为使 $k_{13}$ 最大，只需使 $\frac{m_{2}}{(4m_{0}+m_{2})(m_{2}+m_{0})}=\frac{1}{5m_{0}+m_{2}+4m_{0}^{2}/m_{2}}$ 最大，即 $m_{2}+\frac{4m_{0}^{2}}{m_{2}}$ 取最小值。由 $m_{2}+\frac{4m_{0}^{2}}{m_{2}}=\left(\sqrt{m_{2}}-\frac{2m_{0}}{\sqrt{m_{2}}}\right)^{2}+4m_{0}$ 可知，当 $\sqrt{m_{2}}=\frac{2m_{0}}{\sqrt{m_{2}}}$，即 $m_{2}=2m_{0}$ 时，$k_{13}$ 最大。
}

\end{enumerate}

\end{document}
